\begin{table}[htbp]
\centering
\small
\caption{Problema A (openai/gpt-4o-mini): métrica estimada em zero-shot, few-shot e pelo rótulo real.}
\label{tab:T1_A_gpt4mini}
\begin{tabular}{lrrrrlrrl}
\toprule
 &  &  & \multicolumn{3}{c}{Perceptron} & \multicolumn{3}{c}{NNLS} \\ \cmidrule(lr){4-6} \cmidrule(lr){7-9}
Regime & Acur.\ LLM vs real & Partição LLM (n0/n1) & Fid.\ (vs LLM) & Acur.\ vs real & $\hat{W}$ unit. & Fid.\ (vs LLM) & Acur.\ vs real & $\hat{W}$ unit. \\
\midrule
zero-shot† & 32,4 (28,0–36,7) & 114/36 & 81,3 (78,0–85,3) & 34,7 (17,3–49,3) & [0,09; 0,99] & 84,4 (83,3–86,0) & 23,1 (15,3–27,3) & [0,26; 0,96] \\
4-shot† & 81,9 (67,8–91,8) & 90/56 & 89,2 (84,2–94,5) & 86,1 (68,0–94,0) & [0,18; 0,96] & 89,0 (86,3–94,5) & 89,9 (69,3–96,0) & [0,29; 0,93] \\
10-shot† & 89,4 (81,4–96,4) & 82/58 & 94,4 (90,7–98,6) & 93,4 (86,7–98,0) & [0,34; 0,89] & 93,7 (90,0–97,9) & 94,3 (87,3–98,0) & [0,62; 0,60] \\
20-shot† & 90,9 (84,6–95,4) & 72/58 & 95,5 (92,3–97,7) & 93,9 (89,3–98,0) & [0,42; 0,85] & 95,0 (91,5–97,7) & 94,7 (89,3–98,0) & [0,56; 0,62] \\
40-shot† & 91,4 (83,6–96,4) & 59/51 & 95,2 (90,9–98,2) & 93,4 (88,0–98,0) & [0,50; 0,81] & 93,9 (89,1–98,2) & 94,8 (92,0–98,0) & [0,54; 0,75] \\
\midrule
W ajustado ao rótulo real (teto)† & — & — & — & 96,2 (95,3–98,0) & [0,37; 0,89] & — & 95,3 (93,3–98,0) & [0,51; 0,86] \\
\bottomrule
\multicolumn{9}{p{0.95\linewidth}}{\footnotesize Médias das sementes e repetições, mín–máx entre parênteses, em \%. Fid.\ (vs LLM) = fidelidade: fração dos pontos em que a métrica reproduz o rótulo do LLM. Acur.\ vs real = acurácia (vs rótulo real) da métrica aplicada aos 150 pontos; Acur.\ LLM vs real = acurácia (vs rótulo real) do próprio LLM, mostrada como contexto. Few-shot: exemplos rotulados pelo rótulo real (k/2 por classe, sorteio aleatório), LLM classifica os $150-k$ restantes; fidelidade medida nesses pontos. Teto: W ajustado diretamente ao rótulo real; não tem fidelidade (não é estimado a partir do LLM), só acurácia (vs rótulo real), que não chega a 100\% pela sobreposição das gaussianas do A. Em zero-shot (e com nomes de classe neutros) a correspondência entre nome e classe é arbitrária: a acurácia (vs rótulo real) do LLM pode sair abaixo de 50\%; o valor invariante à permutação é max($a$, 1 $-$ $a$). † = em pelo menos uma estimativa o Perceptron devolveu $\hat{W}$ em melhor esforço ($\gamma = 0$, sem margem viável). ``—'' = colapso em classe única (ou quase-colapso: classe minoritária com menos de 5 pontos); na linha colapsada também a acurácia do LLM e a partição ficam como ``—''.} \\
\end{tabular}
\end{table}
